\documentclass[11pt,a4paper]{article}
\usepackage[utf8]{inputenc}
\usepackage[T1]{fontenc}
\usepackage{lmodern}
\usepackage{amsmath}
\usepackage[margin=2cm]{geometry}
\usepackage{enumitem}
\usepackage{xcolor}
\pagestyle{empty}
\setlength{\parindent}{0pt}
\definecolor{verde}{HTML}{1F8A4F}

\begin{document}

{\LARGE\bfseries Tarea: operaciones combinadas con fracciones}

\vspace{10pt}
Nombre: \rule{9cm}{0.4pt} \hfill Fecha: \rule{3.5cm}{0.4pt}

\vspace{10pt}
{\color{verde}\rule{\textwidth}{1pt}}

\vspace{4pt}
\textbf{Instrucciones.} Resuelve cada ejercicio en tu cuaderno con todo el procedimiento y marca la opción correcta. Simplifica siempre el resultado. Envía la foto de tu procedimiento hoy.

\vspace{10pt}
\begin{enumerate}[label=\textbf{\arabic*.},leftmargin=*,itemsep=14pt]
\item $\displaystyle \frac{7}{12}+\frac{5}{18}-\frac{1}{4}$

\vspace{6pt}
\noindent \textbf{a)}~$\displaystyle \frac{4}{9}$ \hfill \textbf{b)}~$\displaystyle \frac{11}{18}$ \hfill \textbf{c)}~$\displaystyle \frac{13}{18}$ \hfill \textbf{d)}~$\displaystyle \frac{11}{26}$
\vspace{10pt}

\item $\displaystyle \frac{22}{5}-\frac{11}{4}$

\vspace{6pt}
\noindent \textbf{a)}~$\displaystyle 11$ \hfill \textbf{b)}~$\displaystyle \frac{27}{20}$ \hfill \textbf{c)}~$\displaystyle \frac{11}{9}$ \hfill \textbf{d)}~$\displaystyle \frac{33}{20}$
\vspace{10pt}

\item $\displaystyle \frac{5}{6}-\left(\frac{2}{3}-\frac{3}{4}\right)$

\vspace{6pt}
\noindent \textbf{a)}~$\displaystyle \frac{11}{12}$ \hfill \textbf{b)}~$\displaystyle -\frac{7}{12}$ \hfill \textbf{c)}~$\displaystyle \frac{1}{12}$ \hfill \textbf{d)}~$\displaystyle \frac{3}{4}$
\vspace{10pt}

\item $\displaystyle \left(\frac{9}{10}\cdot\frac{25}{27}\right)\div\frac{5}{12}$

\vspace{6pt}
\noindent \textbf{a)}~$\displaystyle \frac{6}{5}$ \hfill \textbf{b)}~$\displaystyle \frac{25}{72}$ \hfill \textbf{c)}~$\displaystyle 2$ \hfill \textbf{d)}~$\displaystyle \frac{1}{2}$
\vspace{10pt}

\item $\displaystyle \frac{3}{4}\div\left(\frac{1}{2}+\frac{1}{3}\right)-\frac{1}{5}$

\vspace{6pt}
\noindent \textbf{a)}~$\displaystyle \frac{49}{30}$ \hfill \textbf{b)}~$\displaystyle \frac{11}{20}$ \hfill \textbf{c)}~$\displaystyle \frac{9}{10}$ \hfill \textbf{d)}~$\displaystyle \frac{7}{10}$
\vspace{10pt}

\item $\displaystyle \frac{2}{3}\left(\frac{3}{2}-\frac{3}{4}\right)+\frac{1}{6}$

\vspace{6pt}
\noindent \textbf{a)}~$\displaystyle \frac{1}{2}$ \hfill \textbf{b)}~$\displaystyle \frac{2}{3}$ \hfill \textbf{c)}~$\displaystyle \frac{7}{12}$ \hfill \textbf{d)}~$\displaystyle \frac{5}{12}$
\vspace{10pt}

\item $\displaystyle \left[\frac{5}{6}+\frac{1}{3}\cdot\frac{3}{4}\right]\div\left[\left(3-\frac{1}{2}\right)\cdot\frac{2}{15}\right]$

\vspace{6pt}
\noindent \textbf{a)}~$\displaystyle \frac{13}{4}$ \hfill \textbf{b)}~$\displaystyle \frac{21}{8}$ \hfill \textbf{c)}~$\displaystyle \frac{13}{36}$ \hfill \textbf{d)}~$\displaystyle \frac{4}{13}$
\vspace{10pt}

\item $\displaystyle \dfrac{2-\dfrac{1}{3}}{1+\dfrac{1}{2+\dfrac{1}{2}}}$

\vspace{6pt}
\noindent \textbf{a)}~$\displaystyle \frac{5}{3}$ \hfill \textbf{b)}~$\displaystyle \frac{21}{25}$ \hfill \textbf{c)}~$\displaystyle \frac{25}{21}$ \hfill \textbf{d)}~$\displaystyle \frac{7}{3}$
\vspace{10pt}

\item $\displaystyle 1-\left[\frac{1}{2}-\left(\frac{1}{3}-\frac{1}{4}\right)\right]$

\vspace{6pt}
\noindent \textbf{a)}~$\displaystyle -\frac{1}{12}$ \hfill \textbf{b)}~$\displaystyle \frac{7}{12}$ \hfill \textbf{c)}~$\displaystyle \frac{13}{12}$ \hfill \textbf{d)}~$\displaystyle \frac{5}{12}$
\vspace{10pt}

\item Un estudiante gasta $\dfrac{1}{4}$ de su dinero en libros y luego $\dfrac{2}{5}$ de lo que le queda en transporte. ¿Qué fracción del dinero total le queda?

\vspace{6pt}
\noindent \textbf{a)}~$\displaystyle \frac{11}{20}$ \hfill \textbf{b)}~$\displaystyle \frac{7}{20}$ \hfill \textbf{c)}~$\displaystyle \frac{3}{10}$ \hfill \textbf{d)}~$\displaystyle \frac{9}{20}$
\vspace{10pt}
\end{enumerate}

\end{document}
